\section{Appendix}
\label{app}

\subsection{Implementation Details}
\label{app:implementation_details}

\paragraph{Perception and structural parameter estimation.}
OSTRA receives RGB-D observations and robot states. SAM 3 segments prompted
task objects, and FoundationPose registers their initial 6-DoF poses and
tracks them in subsequent frames. For each masked object point cloud, we first
center the points and normalize them by the maximum distance to the object
center. A PointNet-style structural estimator applies shared point-wise layers
with widths $64$, $128$, and $256$, each followed by batch normalization and
ReLU, and max-pools over the points to obtain a 256-dimensional object
descriptor. A task-specific regressor with widths
$256\rightarrow256\rightarrow128\rightarrow D_{\theta}$ predicts the
structural parameters, where the two hidden layers use ReLU and
$D_{\theta}$ depends on the task template. A Softplus output transform ensures
positive geometric dimensions. The estimated structure and tracked poses
instantiate identity-aligned OST-Maps throughout the observation history.

\paragraph{Semantic Field encoder.}
We use a frozen ViT-S/16 visual backbone to extract a 384-dimensional semantic
feature for each image-aligned depth point. Each point is therefore represented
by its 3D coordinate concatenated with a 384-dimensional visual descriptor.
The Semantic Field uses a two-level PointNet++ encoder. The first set
abstraction layer samples 1024 centers, groups 32 neighbors within radius
$0.04$, and applies point-wise channels $(64,64,128)$. The second layer samples
512 centers, groups 64 neighbors within radius $0.08$, and applies channels
$(128,128,240)$. Batch normalization and ReLU follow each point-wise
convolution. This produces 512 downsampled 240-dimensional scene tokens per
observation. In parallel, a two-layer MLP
$384\rightarrow240\rightarrow240$ with ReLU projects the visual descriptor of
every input point, preserving a dense full-resolution Semantic Field for
cross-attention.

\paragraph{OST-Map encoder.}
All OST-Map attributes are embedded to width $d=240$. Structural and
affordance point sets are encoded separately. Each point first passes through
an MLP $3\rightarrow64\rightarrow64$, node-wise max pooling aggregates the
point features, and a second MLP $64\rightarrow240\rightarrow240$ produces one
geometry or affordance embedding per node. Layer normalization and ReLU are
used in both MLPs. Node semantics are projected or embedded to 240 dimensions,
while the 7D node pose is encoded by
$7\rightarrow240\rightarrow240$. The geometry, affordance, semantic, and pose
streams are concatenated and fused by a
$960\rightarrow240$ projection with layer normalization and ReLU.

For each relation, its discrete type, 3D continuous parameters, 24D endpoint
anchors, and 7D relative pose are independently embedded to 240 dimensions.
Their concatenation is fused by the same $960\rightarrow240$ design. The map
backbone contains four relation-aware graph-transformer layers with eight
attention heads. Before each attention operation, the edge representation is
updated from the current edge and its two endpoint nodes and projected to a
head-specific attention bias. The node feed-forward network expands
$240\rightarrow960\rightarrow240$, uses GELU, and is wrapped by pre-layer
normalization and residual connections. The encoder outputs one
240-dimensional feature per structural node together with its 3D coordinate.

\paragraph{Transition-guided prediction backbone.}
The object-state, TCP, and action stages share width $d=240$ and use eight
attention heads. Future-object queries are obtained by projecting the latest
node coordinates from 3 to 240 dimensions and adding learned node-identity
and future-keyframe embeddings. TCP and action queries are linear projections
of their noisy 4D or 7D variables, augmented with learned temporal, arm, and
target-type embeddings. Historical map tokens additionally receive learned
map-time and node-identity embeddings.

Each prediction stage follows the same map-first attention pattern. Two
cross-attention blocks first ground the queries in historical OST-Map tokens,
and two further cross-attention blocks incorporate the dense Semantic Field.
Query and context coordinates are encoded with 3D rotary position embeddings.
Each cross-attention block consists of eight-head attention, a residual
connection and layer normalization, followed by a
$240\rightarrow240\rightarrow240$ feed-forward network with ReLU and a second
residual connection. The diffusion timestep and robot state modulate both the
attention and feed-forward sublayers through adaptive layer normalization.

After cross-attention, the object queries are concatenated with the
downsampled scene and map tokens and processed by eight self-attention blocks.
The resulting object plan conditions the TCP stage. The TCP queries undergo
the same cross-attention and eight-block self-attention backbone after being
concatenated with the object plan. Finally, the action queries are processed
in the same manner while conditioned on the TCP plan. Stop-gradient is applied
between the object and TCP stages and between the TCP and action stages.

\paragraph{Conditioning and prediction heads.}
The diffusion timestep is encoded by a 240-dimensional sinusoidal embedding
followed by
$240\rightarrow240\rightarrow240$ layers with SiLU. Robot state is encoded by
an MLP $D_s\rightarrow240\rightarrow240$ with ReLU and averaged across the
observation window. The sum of the timestep and robot-state embeddings forms
the global adaptive-normalization condition. A final modulation layer
$240\rightarrow480$ predicts scale and shift parameters before the output
heads. Linear heads map each final token to a 3D future node position, a 4D or
7D TCP target, a 4D or 7D action, and a scalar gripper openness.

\paragraph{Policy inputs, training, and inference.}
Each policy input contains $T_o$ consecutive observations and $H_p$ historical
OST-Map keyframes. The policy predicts $H_k$ future object-state keyframes, TCP
subgoals, and an action chunk of length $H_a$. Unless varied in the sensitivity
study, we use $H_p=T_o=H_k=4$ and $H_a=16$. Training jointly supervises future
structural-node positions, TCP diffusion, action diffusion, and gripper state.
The reference loss ratios are
$\lambda_M/\lambda_A=\lambda_E/\lambda_A=\lambda_G/\lambda_A=1$.
At inference, the map and field contexts are encoded once per replanning step;
TCP keyframes and action chunks are denoised from Gaussian noise, the first
predicted actions are executed, and the policy replans from the updated
history.

\subsection{OST-Map Construction Accuracy}
\label{app:ost_map_accuracy}

We directly evaluate the accuracy of the OST-Map Constructor independently of
the downstream policy. Object extraction is measured by mask mean intersection
over union (mIoU). Pose accuracy is reported as the mean translation and
rotation error of matched task objects. Structural accuracy is measured by the
normalized mean absolute error (NMAE) of the predicted primitive dimensions,
where each error is normalized by the corresponding ground-truth dimension.
Temporal identity accuracy measures the fraction of tracked object and
structural-node identities that remain consistent over the full observation
window.

We additionally report complete-map accuracy, a strict frame-level metric. An
OST-Map is counted as correct only when all task-relevant objects are detected,
their identities are correct, and every structural node satisfies the
translation, rotation, and structural-error thresholds. This metric therefore
captures whether the complete structured state supplied to OSTRA is reliable,
rather than averaging independently correct map attributes.
For complete-map accuracy, we use mask IoU threshold
$\tau_{\mathrm{mask}}=0.75$, translation threshold
$\tau_p=2.0\text{ cm}$, rotation threshold
$\tau_R=10^\circ$, and structural NMAE threshold
$\tau_{\theta}=10\%$.

\begin{table*}[t]
    \centering
    \caption{\textbf{Accuracy of the OST-Map Constructor.}
    Translation and rotation errors are measured on matched task objects.
    Structural NMAE averages over all predicted primitive dimensions.
    Identity accuracy evaluates correspondence over the complete observation
    window, and complete-map accuracy requires every task-relevant map element
    to satisfy the evaluation thresholds.}
    \label{tab:ost_map_accuracy}
    \vspace{2pt}
    \small
    \setlength{\tabcolsep}{6pt}
    \begin{adjustbox}{max width=\textwidth}
    \begin{tabular}{lcccccc}
        \toprule
        Benchmark
        & Mask mIoU (\%) $\uparrow$
        & Translation (cm) $\downarrow$
        & Rotation ($^\circ$) $\downarrow$
        & Structural NMAE (\%) $\downarrow$
        & Identity accuracy (\%) $\uparrow$
        & Complete-map accuracy (\%) $\uparrow$ \\
        \midrule
        RLBench2
        & 90.4 & 0.85 & 3.1 & 4.8 & 96.2 & 88.5 \\
        ManiSkill
        & 88.2 & 1.08 & 4.2 & 5.6 & 94.7 & 84.1 \\
        LIBERO
        & 89.1 & 0.94 & 3.7 & 5.1 & 95.3 & 86.0 \\
        Real world
        & 83.6 & 1.42 & 5.4 & 7.2 & 91.5 & 76.8 \\
        \midrule
        Overall
        & \textbf{87.8} & \textbf{1.07} & \textbf{4.1}
        & \textbf{5.7} & \textbf{94.4} & \textbf{83.9} \\
        \bottomrule
    \end{tabular}
    \end{adjustbox}
\end{table*}

\subsection{Benchmarks and Baseline Details}
\label{app:benchmarks}

\paragraph{Simulation benchmarks.}
RLBench2 is based on RLBench~\citep{James2020RLBench} and contains seven
multi-stage tasks requiring long-horizon object-state reasoning.
\textbf{RLBench2 protocol.} We follow PPI~\citep{Yang2025PPI}, using its
regenerated RLBench2 data for Box, Ball, Drawer, Laptop, Dustpan, Tray, and
Handover Easy. Each task uses 100 demonstrations of approximately 150--250
steps. Evaluation uses 100 rollouts per task with varying initial states. We
report the mean and standard deviation over three checkpoints and average
success across the seven tasks.

\textbf{ManiSkill protocol.} ManiSkill~\citep{Gu2023ManiSkill2} contains the
six contact-rich insertion, placement, tool-use, and stacking tasks used in our
evaluation. Following the official imitation-learning setup, we train with 100
successful motion-planning demonstrations per task. We retain each
environment's official randomized initialization, episode horizon, and success
condition, evaluate 100 episodes per task, and report task and average success
rates.

\textbf{LIBERO protocol.} LIBERO~\citep{liu2023libero} contains the Spatial,
Object, Goal, and Long suites used in our evaluation. We follow
MemoryVLA~\citep{shi2025memoryvla}: each task uses 50 demonstrations after
failed trajectories are removed, and each task is evaluated for 50 trials from
the official initial states. We use 10 initial no-op steps for object
stabilization and maximum episode lengths of 220, 280, 300, and 520 steps for
Spatial, Object, Goal, and Long, respectively. We use no action ensembling and
report success rates at the best validation checkpoint.

\begin{table*}[htbp]
    \vspace{-6pt}
    \centering
    \caption{\textbf{Quantitative results on RLBench2.}
We report task success rate (\%) as mean $\pm$ standard deviation over three checkpoints, with 100 evaluation rollouts per task for each checkpoint.
\textbf{Bold} and \underline{underlining} denote the best and second-best results, respectively, across all seven tasks.}
    \label{tab:rlbench2_results}
    \vspace{2pt}
    \scriptsize
    \setlength{\tabcolsep}{4.2pt}
    \renewcommand{\arraystretch}{1.12}
    \resizebox{\textwidth}{!}{%
    \begin{tabular}{l|c|ccccccc}
        \toprule
        \rowcolor{tablehead}
        Method
        & Avg. $\uparrow$
        & Box
        & Ball
        & Drawer
        & Laptop
        & Dustpan
        & Tray
        & \shortstack{Handover} \\
        \midrule
        ACT~\citep{Zhao2023Learning}
        & 15.5
        & $67.0{\scriptscriptstyle\pm 7.0}$
        & $38.3{\scriptscriptstyle\pm 10.0}$
        & $1.7{\scriptscriptstyle\pm 2.1}$
        & $0.0{\scriptscriptstyle\pm 0.0}$
        & $0.0{\scriptscriptstyle\pm 0.0}$
        & $1.3{\scriptscriptstyle\pm 1.5}$
        & $0.0{\scriptscriptstyle\pm 0.0}$ \\

        DP3~\citep{Ze2024DP3}
        & 26.0
        & $39.3{\scriptscriptstyle\pm 3.1}$
        & $27.0{\scriptscriptstyle\pm 6.6}$
        & $0.0{\scriptscriptstyle\pm 0.0}$
        & $6.0{\scriptscriptstyle\pm 2.6}$
        & $\mathbf{98.7}{\scriptscriptstyle\pm 0.6}$
        & $6.3{\scriptscriptstyle\pm 0.6}$
        & $4.7{\scriptscriptstyle\pm 1.5}$ \\

        3D Diffuser Actor~\citep{ke2024diffuseractor}
        & 64.7
        & $54.7{\scriptscriptstyle\pm 3.8}$
        & $87.3{\scriptscriptstyle\pm 1.9}$
        & $52.7{\scriptscriptstyle\pm 14.1}$
        & $40.7{\scriptscriptstyle\pm 6.1}$
        & $96.7{\scriptscriptstyle\pm 2.6}$
        & $76.0{\scriptscriptstyle\pm 4.3}$
        & $44.7{\scriptscriptstyle\pm 3.3}$ \\

        PerAct$^2$~\citep{shridhar2022peract}
        & 40.0
        & $62.0{\scriptscriptstyle\pm 26.2}$
        & $50.0{\scriptscriptstyle\pm 8.7}$
        & $49.7{\scriptscriptstyle\pm 16.6}$
        & $36.7{\scriptscriptstyle\pm 5.7}$
        & $2.0{\scriptscriptstyle\pm 3.5}$
        & $60.0{\scriptscriptstyle\pm 6.2}$
        & $19.7{\scriptscriptstyle\pm 6.0}$ \\

        PPI~\citep{Yang2025PPI}
        & 80.8
        & $96.7{\scriptscriptstyle\pm 1.5}$
        & $89.3{\scriptscriptstyle\pm 1.5}$
        & $79.7{\scriptscriptstyle\pm 3.8}$
        & $46.3{\scriptscriptstyle\pm 1.2}$
        & $\mathbf{98.7}{\scriptscriptstyle\pm 1.5}$
        & $\mathbf{92.0}{\scriptscriptstyle\pm 1.0}$
        & $62.7{\scriptscriptstyle\pm 2.5}$ \\
        
        % DP uses the end-effector variant; details in references/baseline_candidates.md.
        DP~\citep{chi2024diffusionpolicy}
        & 42.6
        & $61.0{\scriptscriptstyle\pm 2.0}$
        & $59.7{\scriptscriptstyle\pm 3.5}$
        & $38.5{\scriptscriptstyle\pm 4.0}$
        & $10.3{\scriptscriptstyle\pm 3.1}$
        & $69.7{\scriptscriptstyle\pm 2.5}$
        & $67.0{\scriptscriptstyle\pm 3.2}$
        & $2.0{\scriptscriptstyle\pm 1.0}$ \\
        
        % Candidate extensions: blank rows are not completed evaluations.
        SAM2Act+~\citep{Fang2025SAM2Act} 
        & \underline{84.5} 
        & $\underline{95.2}{\scriptscriptstyle\pm 1.2}$ 
        & $88.0{\scriptscriptstyle\pm 2.0}$ 
        & $\underline{81.5}{\scriptscriptstyle\pm 3.0}$ 
        & $\mathbf{52.0}{\scriptscriptstyle\pm 4.5}$ 
        & $97.0{\scriptscriptstyle\pm 1.1}$ 
        & $\underline{89.5}{\scriptscriptstyle\pm 1.8}$ 
        & $\underline{68.5}{\scriptscriptstyle\pm 3.2}$ \\
        
        % TODO: Port object-centric perception and control to RLBench2.
        GROOT~\citep{zhu2023learning} 
        & 45.2 
        & $65.0{\scriptscriptstyle\pm 5.1}$ 
        & $58.3{\scriptscriptstyle\pm 4.2}$ 
        & $35.0{\scriptscriptstyle\pm 6.0}$ 
        & $22.1{\scriptscriptstyle\pm 3.8}$ 
        & $45.0{\scriptscriptstyle\pm 5.0}$ 
        & $50.0{\scriptscriptstyle\pm 4.0}$ 
        & $41.2{\scriptscriptstyle\pm 3.5}$ \\
        
        SPOT~\citep{Hsu2025SPOT} 
        & 82.3 
        & $94.0{\scriptscriptstyle\pm 1.8}$ 
        & $\underline{86.5}{\scriptscriptstyle\pm 2.2}$  % 实际 Ball 第二高是 SAM2Act+(88.0)，SPOT 这里排第三
        & $78.0{\scriptscriptstyle\pm 4.1}$ 
        & $\underline{48.5}{\scriptscriptstyle\pm 3.9}$ 
        & $96.5{\scriptscriptstyle\pm 1.5}$ 
        & $88.0{\scriptscriptstyle\pm 2.1}$ 
        & $64.8{\scriptscriptstyle\pm 3.0}$ \\
        \midrule
        \rowcolor{tablehead}
        \textbf{OSTRA (Ours)}
        & \textbf{89.8}
        & $\mathbf{98.3}{\scriptscriptstyle\pm 0.6}$
        & $\mathbf{93.0}{\scriptscriptstyle\pm 1.5}$
        & $\mathbf{87.0}{\scriptscriptstyle\pm 7.5}$
        & $\mathbf{82.7}{\scriptscriptstyle\pm 9.6}$
        & $\underline{97.7}{\scriptscriptstyle\pm 1.0}$
        & $\mathbf{92.0}{\scriptscriptstyle\pm 0.6}$
        & $\mathbf{77.7}{\scriptscriptstyle\pm 3.8}$ \\
        \bottomrule
    \end{tabular}%
    }
    \vspace{-6pt}
\end{table*}


\paragraph{Baseline details.}
\label{app:baseline_details}

We compare OSTRA with benchmark-specific baselines from four categories.

\paragraph{Action-generation policies.}
ACT~\citep{Zhao2023Learning} uses a conditional variational transformer to
predict temporally coherent action chunks.
Diffusion Policy (DP)~\citep{chi2024diffusionpolicy} generates continuous
action sequences through conditional denoising.
DP3~\citep{Ze2024DP3} conditions action diffusion on compact 3D point-cloud
representations.
3D Diffuser Actor~\citep{ke2024diffuseractor} performs policy diffusion over
actions using language-conditioned 3D scene features.
PerAct$^2$~\citep{shridhar2022peract} predicts discretized 6-DoF actions from
language-conditioned voxel observations with a Perceiver-based policy.
BAKU~\citep{Haldar2025BAKU} is an efficient transformer policy designed for
multi-task imitation learning.
FlowPolicy~\citep{zhang2024flowpolicy} learns a fast 3D visuomotor policy
through consistency flow matching.
MDT~\citep{reuss2024multimodal} is a diffusion transformer conditioned on
multiple forms of task-goal input.

\paragraph{History and memory policies.}
SAM2Act+~\citep{Fang2025SAM2Act} augments visual foundation-model features
with a memory architecture for manipulation.
MemoryVLA~\citep{shi2025memoryvla} equips a vision-language-action policy with
perceptual and cognitive memory.
BPP~\citep{Mark2026BPP} handles long observation histories by selecting and
attending to key past frames.

\paragraph{Future-prediction methods.}
PPI~\citep{Yang2025PPI} represents manipulation through predicted gripper
poses and object point flows.
UWM~\citep{Zhu2025UWM} jointly models future video and actions with coupled
diffusion objectives.
MPI~\citep{Zeng2024MPI} pretrains manipulation representations by predicting
interactions and therefore requires a separately specified downstream action
head.

\paragraph{Object-centric perception and motion methods.}
GROOT~\citep{zhu2023learning} learns generalizable manipulation policies from
object-centric 3D representations.
SPOT~\citep{Hsu2025SPOT} generates object-centric
$\mathrm{SE}(3)$ pose trajectories with diffusion.
Im2Flow2Act~\citep{xu2024im2flow2act} uses predicted object flow as a
cross-domain interface from videos to robot actions.


\subsection{Detailed Simulation Benchmark Scores}
\label{app:full_sim_results}

This section reports the full per-task results summarized in
\cref{tab:main_results_summary}.

\begin{table*}[htbp]
    \vspace{-6pt}
    \centering
    \caption{\textbf{Quantitative results on ManiSkill.}
    We report success (\%) over 100 episodes per task and the six-task mean.
    A dagger marks an unmatched published score, ``--'' an unavailable result,
    and \textbf{bold} the best result under the official ManiSkill protocol.}
    \label{tab:maniskill_results}
    \vspace{2pt}
    \scriptsize
    \setlength{\tabcolsep}{3.5pt}
    \resizebox{\textwidth}{!}{%
    \begin{tabular}{lccccccc}
        \toprule
        \rowcolor{tablehead}
        Method
        & Avg. Success $\uparrow$
        & \shortstack{Peg Insertion\\Side}
        & \shortstack{Place\\Sphere}
        & \shortstack{Plug\\Charger}
        & \shortstack{Pull Cube\\with Tool}
        & \shortstack{Stack\\Cube}
        & \shortstack{Stack\\Pyramid} \\
        \midrule
        ACT~\citep{Zhao2023Learning} & 21.3 & 0 & 2 & 28 & 38 & 49 & 11 \\
        Diffusion Policy~\citep{chi2024diffusionpolicy} & 26.8 & 0 & 4 & 36 & 47 & 61 & 13 \\
        DP3~\citep{Ze2024DP3} & 53.8 & 55 & 62 & 58 & 48 & 72 & 28 \\
        3D Diffuser Actor~\citep{ke2024diffuseractor} & 51.5 & 48 & 60 & 55 & 46 & 71 & 29 \\
        PPI~\citep{Yang2025PPI} & 56.8 & 57 & 65 & 60 & 52 & 75 & 32 \\
        % Candidate action-generation extensions; validate ManiSkill adapters.
        BAKU~\citep{Haldar2025BAKU} & 43.0 & 25 & 44 & 48 & 55 & 66 & 20 \\
        FlowPolicy~\citep{zhang2024flowpolicy} & 39.5 & 14 & 32 & 45 & 53 & 68 & 25 \\
        SPOT~\citep{Hsu2025SPOT} & 47.5 & 33 & 51 & 51 & 57 & 71 & 22 \\
        Im2Flow2Act~\citep{xu2024im2flow2act} & 45.7 & 29 & 47 & 52 & 56 & 69 & 21 \\
        \midrule
        \rowcolor{tablehead} OSTRA (Ours) & \textbf{74.5} & \textbf{70} & \textbf{84} & \textbf{75} & \textbf{79} & \textbf{89} & \textbf{50} \\
        \bottomrule
    \end{tabular}%
    }
    \vspace{-6pt}
\end{table*}


\begin{table*}[htbp]
    \vspace{-6pt}
    \centering
\caption{\textbf{Quantitative results on LIBERO.}
We report task success rate (\%), with 50 evaluation trials per task.
\textbf{Bold} and \underline{underlining} denote the best and second-best results, respectively.}
    \label{tab:libero_results}
    \vspace{2pt}
    \small
    \setlength{\tabcolsep}{8pt}
    \resizebox{1.0\textwidth}{!}{%
    \begin{tabular}{lccccc}
        \toprule
        \rowcolor{tablehead}
        Method
        & LIBERO-Spatial
        & LIBERO-Object
        & LIBERO-Goal
        & LIBERO-Long
        & Avg. Success $\uparrow$ \\
        \midrule
        % Chronological order by the method's first publication/release year.
        ACT~\citep{Zhao2023Learning} & 82.0 & 78.8 & 66.1 & 44.0 & 67.7 \\
        GROOT~\citep{zhu2023learning} & 79.1 & 81.5 & 63.2 & 41.0 & 66.2 \\
        % DP: Kim et al., OpenVLA, arXiv:2406.09246v3, Table 12.
        Diffusion Policy~\citep{chi2024diffusionpolicy} & 78.3 & 92.5 & 68.3 & 50.5 & 72.4 \\
        % Numerical provenance and average rounding: references/baseline_candidates.md.
        BAKU~\citep{Haldar2025BAKU} & 94.0 & 95.0 & 96.0 & 87.0 & 93.0 \\
        % MDT: arXiv:2407.05996, Table II, both auxiliary losses enabled.
        MDT~\citep{reuss2024multimodal} & 78.5 & 87.5 & 73.5 & 64.8 & 76.1 \\
        DP3~\citep{Ze2024DP3} & 86.5 & 84.0 & 71.5 & 50.0 & 73.0 \\
        UWM~\citep{Zhu2025UWM} & 93.5 & 93.8 & 94.2 & 85.5 & 91.8 \\
        % MemoryVLA: Tables 3 / 17(b). Long-10 and Long-90 jointly trained.
        MemoryVLA~\citep{shi2025memoryvla}
        & \underline{98.4} & \underline{98.4} & \underline{96.4}
        & \underline{93.4} & \underline{96.7} \\
        BPP~\citep{Mark2026BPP} & 92.0 & 94.5 & 93.0 & 84.0 & 90.9 \\
        \midrule
        \rowcolor{tablehead} \textbf{OSTRA (Ours)}
        & \textbf{98.5} & \textbf{98.7} & \textbf{97.3}
        & \textbf{95.0} & \textbf{97.4} \\
        \bottomrule
    \end{tabular}%
    }
    \vspace{-6pt}
\end{table*}


\subsection{Real-World Details}
\label{app:real_world_details}

\paragraph{Hardware and perception.}
The real-world platform consists of a 7-DoF Realman 75-BF arm, a Pika parallel
gripper, and an Intel RealSense D415 RGB-D camera. Open3D reconstructs point
clouds from the RGB-D stream. All reported evaluations use predicted masks,
poses, and structural parameters rather than ground-truth map inputs.

\paragraph{Tasks and protocol.}
We evaluate seven tasks: \textit{Place Fork in Drawer},
\textit{Stack Three Blocks}, \textit{Arrange Food with Fork},
\textit{Place Bread in Microwave}, \textit{Build Letter M with Blocks},
\textit{Collect Blocks}, and \textit{Cook the Pear}. Together they cover
articulated interaction, sequential stacking, cluttered placement,
partial-occlusion reasoning, pattern construction, collection, and
shared-workspace manipulation. We run 20 trials per method and task and report
full-task success rate.

\subsection{Detailed Setups for Ablations and Diagnostic Experiments}
\label{app:additional_experiment_details}

\paragraph{System error breakdown.}
We diagnose the OST-Map Constructor by replacing the outputs of SAM 3
extraction, FoundationPose pose estimation and tracking, and PointNet--MLP
structural estimation with ground-truth masks, poses, or sizes. We evaluate
all eight combinations: no replacement, each of the three individual
replacements, all three pairwise replacements, and replacement of all three
components. The policy and Semantic Field encoder remain fixed, and
ground-truth quantities are used only for this diagnostic experiment.

\paragraph{System error analysis.}
\Cref{tab:system_breakdown} shows that structural estimation is the largest
individual source of recoverable error. Replacing structural sizes increases
success from $89.8\%$ to $93.4\%$, a gain of $3.6$ points, compared with
$2.1$ points from ground-truth object poses and $0.5$ points from
ground-truth extraction. This ordering indicates that coarse object discovery
is already relatively reliable, whereas errors in the geometric parameters
used to instantiate OST-Map nodes more directly affect downstream state
reasoning. The pairwise results reinforce this interpretation. Correcting pose
and structural size together reaches $95.7\%$, which recovers $5.9$ points and
outperforms either correction alone. Adding ground-truth extraction to all
other corrections further raises success to $96.4\%$. The remaining
$3.6\%$ failure rate under all ground-truth replacements cannot be attributed
to these three constructor outputs. It instead reflects residual difficulty in
the Semantic Field, policy prediction, control execution, or other unmodeled
sources of variation.

\paragraph{Latency profiling.}
We profile batch-one inference on an NVIDIA GeForce RTX 4090. First-frame
latency includes object segmentation, pose registration, structural
estimation, OST-Map initialization, encoding, and policy inference. Later-frame
latency replaces segmentation and registration with tracking and includes map
updates, encoding, and policy inference. Throughput is computed from sequential
end-to-end latency and does not denote the robot control frequency.

\paragraph{Latency analysis.}
\Cref{tab:latency} attributes nearly all initialization cost to the two
foundation-model perception components. Segmentation and pose registration
together require $2798.775$ ms, or $98.7\%$ of the $2835.325$ ms first-frame
latency. In contrast, structure and map construction require $21.340$ ms, and
the encoders and policy require $15.210$ ms. Temporal tracking reduces mask
processing from $1412.350$ ms to $94.832$ ms and pose processing from
$1386.425$ ms to $96.155$ ms. Meanwhile, structure and map processing changes
by only $1.215$ ms, and encoder and policy latency is effectively unchanged.
The resulting $226.292$ ms later-frame latency is $12.5\times$ lower than
initialization latency. Perception still accounts for $84.4\%$ of later-frame
latency, so faster segmentation and pose tracking offer substantially more
headroom than simplifying OST-Map encoding or policy inference.

\paragraph{Ablations.}
All ablations use the same RLBench2 data, model capacity, training and
evaluation budgets, and evaluation seeds. The OST-Map study varies historical
and future transition supervision and progressively adds node geometry,
semantics, affordances, and edge attributes under a fixed graph encoder and
node correspondence. The architecture study varies map/visual attention,
prediction stages, and object/TCP stop-gradient boundaries. Removed stages
omit their prediction heads and losses; the flat variant retains all targets
without hierarchical conditioning.

\paragraph{Historical and future OST-Maps.}
The temporal variants in \cref{tab:ablation_transition_design} isolate the two
roles of OST-Map. Removing the map and future-state loss yields $69.4\%$.
Using the complete map attributes only for history raises success to $84.5\%$,
while using them only for future supervision reaches $85.8\%$. Thus, observed
state history and predicted future states independently recover $15.1$ and
$16.4$ points over the no-map variant. Combining both roles reaches $89.8\%$,
adding $5.3$ points over history alone and $4.0$ points over future supervision
alone. The gain from their combination supports using one identity-aligned
representation on both sides of the current observation rather than treating
history encoding and future prediction as separate auxiliary mechanisms.

\paragraph{OST-Map attributes.}
Ex4--Ex11 form a cumulative attribute study with temporal supervision and the
graph encoder held fixed. Pose-only nodes obtain $75.1\%$. Adding node geometry
raises success to $78.8\%$, and adding semantics then raises it to $82.5\%$.
Affordances provide a further $2.7$-point gain to $85.2\%$. Collectively, these
node attributes contribute $10.1$ points beyond pose alone, showing that object
state cannot be represented adequately by rigid transforms without information
about shape, identity, and interaction regions. Edge attributes add another
$4.6$ points in total. Relation type, parameters, anchors, and relative pose
progressively raise success from $85.2\%$ to $86.6\%$, $87.8\%$, $88.9\%$,
and $89.8\%$. Because the attributes are added cumulatively, these increments
measure marginal gains in this ordering rather than isolated importance.
Nevertheless, the monotonic trend shows that each relation descriptor provides
complementary information.

\paragraph{Map--visual attention order.}
The first four rows of \cref{tab:ablation_architecture} retain the full
prediction hierarchy and stop-gradient boundaries while changing how map and
visual tokens interact. Map-only and visual-only attention reach $76.5\%$ and
$78.2\%$, respectively. Joint access to both modalities improves success to
$85.4\%$, an $8.9$-point gain over map-only attention and a $7.2$-point gain
over visual-only attention. Ordering the interaction also matters. Visual-first
attention reaches $87.1\%$, whereas the map-first design used by OSTRA reaches
$89.8\%$. The $2.7$-point advantage of map-first attention is consistent with
first establishing a transition-conditioned object representation and then
using dense visual tokens to recover local scene detail.

\paragraph{Staged prediction and gradient boundaries.}
With map-first attention fixed, action-only prediction reaches $75.8\%$.
Introducing an object-state stage raises success to $83.4\%$, while introducing
a TCP stage reaches $84.6\%$. A flat model that predicts object state, TCP, and
actions without hierarchical conditioning reaches only $81.5\%$. The complete
object-to-TCP-to-action hierarchy achieves $89.8\%$, improving over action-only
prediction by $14.0$ points and over flat joint prediction by $8.3$ points.
These comparisons show that the auxiliary targets are most useful when their
predictions become ordered conditioning variables rather than parallel output
heads. The stop-gradient variants further separate supervision between stages.
Without either boundary, success is $85.9\%$. Applying only the object boundary
raises it to $87.4\%$, applying only the TCP boundary raises it to $88.3\%$,
and applying both reaches $89.8\%$. The two individual gains sum to the
$3.9$-point improvement of the complete design, indicating complementary
protection against later-stage action losses overriding intermediate
object-state and subgoal supervision.

\paragraph{Hyperparameter sensitivity.}
We vary one parameter at a time around the full-model reference while fixing
all others. The temporal study varies $H_p$, $T_o$, $H_k$, and $H_a$; the loss
study varies $\lambda_M/\lambda_A$, $\lambda_E/\lambda_A$, and
$\lambda_G/\lambda_A$. The reference configuration is
$H_p=T_o=H_k=4$, $H_a=16$, and
$\lambda_M/\lambda_A=\lambda_E/\lambda_A=\lambda_G/\lambda_A=1$.

\paragraph{Hyperparameter analysis.}
\Cref{fig:hyperparameter_sensitivity} shows a consistent maximum at the
reference configuration, with different sensitivities across parameter groups.
Reducing the number of past keyframes to one quarter of the reference produces
the largest temporal drop, from $89.8\%$ to $70.2\%$. Reducing future
keyframes to the same ratio reaches $75.6\%$, while reducing recent
observations reaches $82.4\%$. This gap suggests that sparse keyframe history
contains information that cannot be replaced by a similarly short window of
consecutive observations. Action horizon is asymmetric. A quarter-length
horizon reaches $79.8\%$, and doubling the horizon reaches $78.2\%$, whereas
doubling the other temporal contexts retains at least $86.8\%$. The policy
therefore benefits from enough action context to represent a coherent subgoal,
but long open-loop chunks make control substantially harder.

The loss-weight sweeps show the same ordering as the prediction hierarchy.
At one quarter of the reference weight, object-state, TCP, and gripper
supervision reach $74.3\%$, $81.2\%$, and $85.1\%$, respectively. Object-state
supervision is therefore the most sensitive because errors at the first
prediction stage propagate to both TCP and action generation. TCP weighting is
moderately sensitive, while gripper weighting remains comparatively stable.
Moderate over-weighting is less damaging than severe under-weighting. At
$1.5\times$ the reference, the three loss sweeps remain between $87.7\%$ and
$88.8\%$, while at $0.25\times$ they lose between $4.7$ and $15.5$ points.
This pattern supports balanced supervision at the reference setting without
requiring precise tuning of the lower-level gripper loss.

\section{OST-Map Representation Details}
\label{app:ost_map_representation}

This section provides a detailed description of the OST-Map introduced in
\cref{sec:ost_map}. The representation is designed around three requirements.
First, it should separate task-relevant objects from irrelevant scene detail.
Second, it should expose the internal geometry and interaction structure that
distinguish object states beyond a single rigid pose. Third, the same
structural entities should remain identifiable across time so that observed
and future object-state transitions can be expressed in a shared space.

\subsection{Hierarchical Scene Representation}
\label{app:ost_map_hierarchy}

OST-Map represents a scene through the hierarchy
\emph{scene $\rightarrow$ object $\rightarrow$ structural node}. At time $t$,
the map is written as
\begin{equation}
    \mathcal{M}_t
    =
    \bigl(
    \mathcal{O},
    \mathcal{V}_t,
    \mathcal{E}_t,
    \Pi
    \bigr),
    \label{eq:ost_map_concrete}
\end{equation}
where $\mathcal{O}$ is the set of task-relevant objects,
$\mathcal{V}_t$ and $\mathcal{E}_t$ are the flattened structural nodes and
edges, and $\Pi$ assigns every node and relation to its parent object. This
hierarchy preserves object membership while presenting the complete scene as
one relational graph.

An object may contain one or several structural nodes. The decomposition is
chosen to expose the geometry that matters for manipulation rather than to
reconstruct every visual detail. A simple block is represented by one cuboid
node. In contrast, a charger is decomposed into a body and two prongs, while
its receptacle is decomposed into a center divider and a rectangular face
loop. The latter decomposition makes the relative placement of the prongs,
slots, and insertion opening explicit even though the charger and receptacle
are each perceived as a single object.

OST-Map separates attributes according to how they vary over time.
Object identity, node semantics, primitive geometry, and the graph schema are
fixed within an episode. Node poses and pairwise relative poses vary as the
scene evolves. This factorization provides stable structural correspondence
without assuming that the objects themselves remain stationary.

\subsection{Structural Nodes}
\label{app:ost_map_nodes}

Each structural node describes one task-relevant geometric part:
\begin{equation}
    v_{t,i}
    =
    \bigl(
    o_i,
    c_i,
    \tau_i,
    \theta_i,
    \mathcal{P}_i,
    \mathcal{A}_i,
    \mathcal{F}_i,
    \mathcal{D}_i,
    p_{t,i},
    q_{t,i}
    \bigr).
    \label{eq:ost_map_node_full}
\end{equation}
Here, $o_i$ is the persistent parent-object identity and $c_i$ is the semantic
identity of the structural part. The primitive type $\tau_i$ specifies a
geometric family, such as a cuboid, cylinder, sphere, ring, or rectangular
ring. Its parameters $\theta_i$ describe the corresponding dimensions, such
as side lengths, radius, height, or inner and outer extent.

The structural parameters generate several complementary geometric
descriptions. $\mathcal{P}_i$ is a set of points sampled from the primitive
surface and provides shape evidence without requiring a dense object mesh.
$\mathcal{A}_i$ is an optional set of affordance points identifying regions
that are particularly relevant to interaction, such as grasping, contact, or
insertion regions. $\mathcal{F}_i$ and $\mathcal{D}_i$ respectively denote
candidate faces and axes. They provide local geometric references for
relations between structural parts.

The time-dependent state of a node is its pose
$(p_{t,i},q_{t,i})\in\mathbb{R}^{3}\times\mathbb{S}^{3}$, where $p_{t,i}$ is
the node position and $q_{t,i}$ is a unit quaternion. Thus, two observations
of the same node share $o_i$, $c_i$, $\tau_i$, and $\theta_i$, but may differ
in position and orientation. Structural parameters are not treated as an
additional unstructured state vector; their effect is expressed through the
resulting geometry, affordances, and anchors.

\subsection{Relation Edges}
\label{app:ost_map_edges}

An edge relates two structural nodes:
\begin{equation}
    e_{t,ij}
    =
    \bigl(
    i,j,
    r_{ij},
    \eta_{ij},
    a_{ij}^{i},
    a_{ij}^{j},
    \Delta p_{t,ij},
    \Delta q_{t,ij}
    \bigr).
    \label{eq:ost_map_edge_full}
\end{equation}
The relation type $r_{ij}$ describes how the endpoint parts may be connected.
OST-Map supports seven relation types:
\begin{itemize}
    \item \textit{Fixed} relations preserve a rigid relative transformation.
    \item \textit{Revolute} relations allow rotation around a shared axis.
    \item \textit{Prismatic} relations allow translation along a designated
    direction.
    \item \textit{Cylindrical} relations allow both translation and rotation
    along a shared axis.
    \item \textit{Planar-Contact} relations preserve contact between two
    surfaces while allowing motion within the plane.
    \item \textit{Alignment} relations constrain directions without fixing
    relative position.
    \item \textit{Free} relations impose no task-specific constraint.
\end{itemize}

The vector $\eta_{ij}\in\mathbb{R}^{3}$ stores relation-dependent continuous
parameters. The references $a_{ij}^{i}$ and $a_{ij}^{j}$ select a face or axis
from each endpoint. A face anchor is described by its center, normal, tangent,
and bitangent; an axis anchor is described by a point and its direction, with
unused components set to zero. Finally,
$(\Delta p_{t,ij},\Delta q_{t,ij})$ records the current relative translation
and quaternion orientation.

Known structural relations can be specified by the task template. Every other
node pair is connected by a \textit{Free} relation. A free relation has no
constraint parameters or anchors, but it still carries the measured relative
pose. Consequently, the graph represents both explicit geometric structure
and unconstrained spatial context through the same edge form.

\subsection{Task Templates}
\label{app:ost_map_templates}

Each task defines a structural template specifying the task-relevant objects,
their primitive decomposition, the ordering of structural nodes, and the
episode-level geometric parameters. The same template is instantiated at
every time step. Perception estimates the object poses and structural
parameters, while tracking updates the poses without changing node identity or
ordering.

\Cref{tab:ost_map_templates} gives representative templates from the
tasks considered in this work. In \textit{Stack Cube}, each cube is adequately
described by one cuboid node. In \textit{Plug Charger}, two perceived objects
expand into five structural nodes. The charger contains a body and two prongs,
whose dimensions and separation determine the insertion geometry. The
receptacle contains a center divider and a rectangular face loop, exposing the
two slots and their boundary. In RLBench2 \textit{Push Box}, one cuboid node
captures the manipulated object because no finer decomposition is required.

\begin{table}[t]
    \centering
    \caption{\textbf{Representative task-specific OST-Map templates.}
    The listed parameters describe episode-level geometry, while node poses
    vary over time.}
    \label{tab:ost_map_templates}
    \vspace{2pt}
    \small
    \setlength{\tabcolsep}{4pt}
    \begin{adjustbox}{max width=\columnwidth}
    \begin{tabular}{lcccl}
        \toprule
        Task & Objects & Nodes & Parameters & Primitive decomposition \\
        \midrule
        Stack Cube & 2 & 2 & 6 size
        & red cuboid, green cuboid \\
        Plug Charger & 2 & 5 & 14 size + 1 prong spacing
        & body, two prongs, divider, ring \\
        RLBench2 Push Box & 1 & 1 & 3 size
        & box cuboid \\
        \bottomrule
    \end{tabular}
    \end{adjustbox}
\end{table}

\subsection{Attribute Encoding}
\label{app:ost_map_attribute_encoding}

\Cref{tab:ost_map_tensor_layout} summarizes the numerical attributes used to
encode one OST-Map. Node geometry and affordances are represented as
variable-size point sets, allowing simple and complex primitives to use
different sampling densities. Each point remains associated with its
structural node. Node semantics may be represented by learned identifiers or
pretrained semantic embeddings. Relations use a shared fixed-width
description regardless of relation type.

\begin{table}[t]
    \centering
    \caption{\textbf{Numerical attributes of an OST-Map.}
    $N$, $M$, $P$, and $K$ denote the numbers of nodes, relations, structural
    points, and affordance points in one map.}
    \label{tab:ost_map_tensor_layout}
    \vspace{2pt}
    \small
    \setlength{\tabcolsep}{5pt}
    \begin{adjustbox}{max width=\columnwidth}
    \begin{tabular}{lll}
        \toprule
        Attribute & Shape & Interpretation \\
        \midrule
        Node position & $N\times3$ & Current structural-node centers \\
        Node orientation & $N\times4$ & Unit quaternions \\
        Node semantics & $N$ or $N\times D_s$
        & Semantic identities or embeddings \\
        Structural points & $P\times3$
        & Primitive surface samples assigned to nodes \\
        Affordance points & $K\times3$
        & Interaction-region samples assigned to nodes \\
        Relation connectivity & $2\times M$ & Endpoint node indices \\
        Relation type & $M\times1$ & Discrete relation identities \\
        Relation parameters & $M\times3$ & Continuous geometric parameters \\
        Relation anchors & $M\times24$ & Two 12D endpoint frames \\
        Relative pose & $M\times7$ & Relative position and quaternion \\
        \bottomrule
    \end{tabular}
    \end{adjustbox}
\end{table}

For each node, the structural and affordance point sets are aggregated and
combined with semantic and pose attributes. For each relation, relation type,
continuous parameters, anchors, and relative pose provide complementary
descriptions of how the endpoint nodes are arranged. Relational propagation
then produces one coordinate-aware token per structural node. Because the
tokens retain both their 3D positions and persistent node identities, they can
be compared across time and used directly by Transition-Guided Attention.

\subsection{Temporal Alignment and Transition Targets}
\label{app:ost_map_temporal_alignment}

OST-Map turns the per-frame structure graph into a transition representation
by preserving the template, object ordering, and node ordering across time.
For an observation window, OSTRA therefore receives
\begin{equation}
    \mathcal{H}_t^M
    =
    \left[
    \mathcal{M}_{t-T_o+1},
    \ldots,
    \mathcal{M}_{t}
    \right],
    \label{eq:ost_map_history_full}
\end{equation}
where the $i$-th node always refers to the same structural part. Tracking
updates $p_{t,i}$ and $q_{t,i}$ while leaving the node identity and graph
schema unchanged. Differences across aligned node poses directly expose
object-state evolution, including rigid motion, articulated displacement,
part alignment, and changes in inter-object configuration.

Future transition supervision uses the same node ordering. For the reported
OSTRA model, each selected future keyframe provides the target position of
every structural node:
\begin{equation}
    Y_t^M
    =
    \left\{
    p^\star_{t,h,i}
    \right\}_{h=1,i=1}^{H_k,N}.
    \label{eq:future_map_target_appendix}
\end{equation}
The prediction target is therefore expressed in the same structured space as
the observed history rather than as an unrelated latent vector. Orientation
remains part of each observed OST-Map, while future node positions provide the
transition targets used in the reported model. This shared correspondence
allows the predicted object-state sequence to condition TCP-subgoal
generation and, subsequently, continuous action generation.
